% Generated draft: requires mathematical and source review.
\documentclass[11pt]{article}
\usepackage[margin=25mm]{geometry}
\usepackage{fontspec}
\setmainfont{Latin Modern Roman}
\usepackage{amsmath,amssymb,mathtools,unicode-math}
\setmathfont{Latin Modern Math}
\usepackage{enumitem,needspace,xurl,hyperref}
\setlength{\parindent}{0pt}
\setlength{\parskip}{5pt}
\setlength{\emergencystretch}{3em}
\newcommand{\N}{\mathbb{N}}
\newcommand{\Z}{\mathbb{Z}}
\newcommand{\Q}{\mathbb{Q}}
\newcommand{\R}{\mathbb{R}}
\newcommand{\C}{\mathbb{C}}
\newcommand{\F}{\mathbb{F}}
\newcommand{\A}{\mathbb{A}}
\newcommand{\PP}{\mathbb{P}}
\newcommand{\E}{\mathbb{E}}
\newcommand{\eps}{\varepsilon}
\newcommand{\dd}{\,\mathrm{d}}
\begin{document}
\section*{UnsolvedMath ID: 30003765}
\textbf{Source problem code:} catalogue entry
\subsection{Definitions and mathematical statement}
Let \(k/\Q\) be a finitely generated field extension. Fix a separable closure \(\bar k\) and write \(G_k=\operatorname{Gal}(\bar k/k)\), with its profinite topology. For a geometrically connected smooth \(k\)-variety \(X\) with geometric base point \(\bar x\), its arithmetic étale fundamental group fits into the exact sequence
\[
1\longrightarrow \pi_1^{\acute{e}t}(X_{\bar k},\bar x)
\longrightarrow \pi_1^{\acute{e}t}(X,\bar x)
\xrightarrow{p_X}G_k\longrightarrow 1.
\]
Base points and their associated inner-conjugacy indeterminacy must be accounted for when comparing these groups. The projection \(p_X\) is additional structure; specifying the two abstract groups alone does not specify this extension.

A smooth geometrically connected curve \(C/k\) is hyperbolic if its smooth projective compactification \(\overline C\) has geometric genus \(g\), its geometric boundary \((\overline C\setminus C)(\bar k)\) has cardinality \(r\), and
\[
2g-2+r>0.
\]
Here \(r=0\) is allowed.

\textbf{Recorded questions and formulation qualifications.} The source asks the following package of reconstruction questions over every such field \(k\).
\begin{enumerate}
\item Are smooth hyperbolic curves over \(k\) anabelian? A usual isomorphism formulation asks whether, for every pair of smooth hyperbolic curves \(C,D/k\), the natural map
\[
\operatorname{Isom}_k(C,D)
\longrightarrow
\operatorname{Isom}^{\mathrm{out}}_{G_k}
\bigl(\pi_1^{\acute{e}t}(C),\pi_1^{\acute{e}t}(D)\bigr)
\]
is bijective. On the right, one takes continuous group isomorphisms \(\varphi\) satisfying \(p_D\circ\varphi=p_C\), modulo conjugation by elements of \(\pi_1^{\acute{e}t}(D_{\bar k})\). This specifies a conventional isomorphism version of the question. The supplied record does not specify whether the word ``anabelian'' also requires a reconstruction statement for a designated class of morphisms.
\item Does every point of every smooth \(k\)-variety \(X\) have a basis of anabelian Zariski neighborhoods? The basis condition means that, for each point \(x\in X\) and each Zariski open neighborhood \(V\subseteq X\) of \(x\), there is a Zariski open subset \(U\) such that
\[
x\in U\subseteq V
\quad\text{and}\quad U\text{ is anabelian over }k.
\]
The source does not define anabelianity for these higher-dimensional neighborhoods, its comparison category, or its required morphisms. It also does not clarify whether ``every point'' means every scheme point or only a specified class of points. The displayed quantification makes the literal scheme-point reading explicit without resolving that uncertainty.
\item Can the set \(X(k)=\operatorname{Hom}_k(\operatorname{Spec}k,X)\) be reconstructed from \(\pi_1^{\acute{e}t}(X)\) together with \(G_k\)? The source does not specify the precise class of \(X\) in this clause, what structure accompanies the two groups, or the meaning and functoriality of reconstruction. When \(X\) is geometrically connected and the projection \(p_X\) is included, each rational point induces a conjugacy class of continuous sections of \(p_X\). Thus there is a natural map
\[
X(k)\longrightarrow
\left\{s:G_k\longrightarrow\pi_1^{\acute{e}t}(X)\ \middle|\
s\text{ is continuous and }p_X\circ s=\operatorname{id}_{G_k}\right\}
\big/\pi_1^{\acute{e}t}(X_{\bar k}),
\]
where the quotient is by conjugation. This describes one possible reconstruction mechanism; the supplied record does not assert that this map is the intended target or provide hypotheses under which it should be bijective. In particular, no such bijectivity assertion for all smooth varieties is supplied here.
\end{enumerate}
For varieties that are not geometrically connected, the single exact sequence above does not supply a formulation without further component or groupoid data. The source does not resolve this issue. These questions therefore retain explicit unresolved qualifications and do not constitute one fully specified conjecture in the recorded generality.
\subsection{Short English statement}
Over every field finitely generated over \(\Q\), the record asks whether hyperbolic curves are determined by their arithmetic fundamental groups, whether smooth varieties have arbitrarily small anabelian Zariski neighborhoods, and whether rational points can be recovered from fundamental-group data. The meaning of higher-dimensional anabelianity and the scope and mechanism of rational-point reconstruction remain unspecified. The supplied status is partially solved; this draft does not independently verify that assessment.
\subsection{Sources}
\begin{itemize}
\item Local UnsolvedMath registry: \texttt{data/unsolvedmath/problems.json}, supplied record with permanent ID \(30003765\), source problem code \texttt{catalogue entry}, titled ``Grothendieck's Anabelian Conjectures over Finitely Generated Fields.'' The supplied record labels the statement \texttt{corrected\_verified} and the problem \texttt{partially\_solved}; these are registry labels, not independent mathematical verification. No individual proposer is identified; the proposed year is recorded as 2018.
\item Supplied report citation: \emph{Mini-Workshop: Arithmetic Geometry and Symmetries around Galois and Fundamental Groups} (2019), Oberwolfach Reports, pp.~1021--1080, DOI \texttt{10.4171/owr/2018/17}. \url{https://doi.org/10.4171/owr/2018/17}. The background attributes the collection to workshop reports and their open-problem discussions. This same URL is listed among the supplied literature sources as ``OWR 2018 report.'' The differing year indications are retained without reconciliation.
\item Supplied literature source: Schmidt--Stix, ``anabelian varieties,'' as labeled in the record. \url{https://doi.org/10.4007/annals.2016.184.3.5}. No fuller bibliographic title or theorem locator is supplied.
\item Literature assessment checked at \texttt{2026-08-21}: this is supplied dated curation, not independently verified. It assigns the status ``partially solved'' and states that Schmidt--Stix prove an arbitrary-dimensional neighborhood theorem and many higher-dimensional cases, while the full package of Grothendieck's reconstruction conjectures in the recorded generality is not proved. The background describes this as dated literature triage derived from the supplied OWR curation and expressly calls for independent verification of time-sensitive status and cited claims. No source was visited for this draft.
\end{itemize}
\end{document}
